% Source: https://artofproblemsolving.com/community/c3h2792393_2016_math_hour_olympiad__university_of_washington__grades_67
% Stable identity: aops-topic-2792393-q01
A witch passes exactly when all $10$ of her non-neighbours fail.
\begin{itemize}
\item Two passing witches must be neighbours, since otherwise each predicted the other would fail. In a circle of $13$, at most $2$ witches are pairwise adjacent, so at most $2$ pass.
\item If nobody passes, every prediction is correct, so everybody passes: contradiction.
\item If exactly one witch $W$ passes, consider a neighbour $N$ of $W$. Every non-neighbour of $N$ fails ($W$ is $N$'s neighbour), so $N$'s predictions are all correct and $N$ passes: contradiction.
\item Exactly two adjacent witches passing is consistent: each of them has only failing non-neighbours, and every other witch has at least one of the two among her non-neighbours (no witch neighbours both), so she fails.
\end{itemize}
So exactly $\boxed{2}$ witches pass.
